\SourceLeaf{157}
\Needspace{6\baselineskip}\pdfbookmark[3]{II. Evaluirung der Verbindlichkeiten gegen gegenwärtige Theilnehmer}{witwen-II-2-II}
\par\medskip{\raggedright\noindent II.\quad Evaluirung des jetzigen Geldwerths der Verbindlichkeiten der Witwenkasse gegen ihre gegenwärtigen Theilnehmer.\par}\smallskip
Wie die Witwen, so habe ich auch alle Interessenten der Witwenkasse nach der Reihenfolge ihres Beitritts mit fortlaufender Numerirung bezeichnet. Indem ich ganz den a. a. O. vorgezeichneten Weg verfolge, habe ich zuvörderst die 42 verheiratheten Theilnehmer in Betracht zu ziehen, und zunächst (was gleichsam den Kern der Untersuchung bildet) den gegenwärtigen Geldwerth theils von den Beiträgen, zu welchen sie verpflichtet sind, theils der Pensionen welche ihre dermaligen Ehefrauen im Fall des Überlebens zu geniessen haben werden, nach den Principien der Wahrscheinlichkeitsrechnung und den Mortalitätstafeln zu ermitteln. Ich concentrire hier auf Einer Seite die Resultate einer langen Arbeit in einem tabellarischen Abrisse.

\begin{center}
\scriptsize\setlength{\tabcolsep}{1pt}
\begin{tabular}{r|c|r|r|r|r||r|c|r|r|r|r}
&\shortstack{Jetzige\\verhei-\\rathete\\Theil-\\nehm.}&\multicolumn{2}{c|}{Zinsfuss \(3\tfrac{1}{2}\) Proc.}&\multicolumn{2}{c||}{Zinsfuss 4 Proc.}&&\shortstack{Jetzige\\verhei-\\rathete\\Theil-\\nehm.}&\multicolumn{2}{c|}{Zinsfuss \(3\tfrac{1}{2}\) Proc.}&\multicolumn{2}{c}{Zinsfuss 4 Proc.}\\\cline{3-6}\cline{9-12}
Nro.&&\multicolumn{2}{c|}{\shortstack{Jetziger Geldwerth\\der}}&\multicolumn{2}{c||}{\shortstack{Jetziger Geldwerth\\der}}&Nro.&&\multicolumn{2}{c|}{\shortstack{Jetziger Geldwerth\\der}}&\multicolumn{2}{c}{\shortstack{Jetziger Geldwerth\\der}}\\
&&Pensionen&Beiträge&Pensionen&Beiträge&&&Pensionen&Beiträge&Pensionen&Beiträge\\\hline
&&Thl.&Thl.&Thl.&Thl.&&&Thl.&Thl.&Thl.&Thl.\\
109&L.&1752&59&1634&58&174&H.&1175&120&1053&114\\
135&O.&1509&85&1386&83&175&R.&1154&142&1007&134\\
138&C.&1846&70&1706&68&177&W.&1117&131&991&124\\
139&U.&937&109&854&104&180&T.&1064&132&944&125\\
140&H.&1451&106&1310&101&183&V.&943&148&824&139\\
141&L.&1091&93&1004&90&184&B.&1023&146&893&137\\
147&R.&1441&105&1301&101&187&W.&917&134&817&127\\
149&G.&1717&101&1545&97&188&S.&867&134&774&127\\
151&O.&1051&130&933&123&189&H.&1138&129&1011&122\\
153&B.&865&109&790&104&190&H.&936&142&825&134\\
155&K.&925&114&840&109&191&D.&892&123&804&117\\
156&S.&853&114&776&109&192&R.&943&139&848&131\\
160&H.&1376&123&1221&117&193&G.&1047&148&907&140\\
161&B.&876&120&795&114&195&D.&981&140&866&132\\
164&W.&1359&122&1209&117&197&R.&916&154&791&145\\
167&S.&995&140&878&132&198&W.&1072&144&936&136\\
168&K.&1345&135&1173&128&199&F.&1127&127&1005&120\\
169&Z.&1118&134&988&127&200&L.&861&153&749&144\\
170&R.&1541&99&1397&96&201&W.&1073&145&935&137\\
171&C.&1193&105&1084&101&203&B.&932&127&835&121\\
173&F.&1120&127&998&121&204&E.&787&146&729&138\\\cline{6-12}
\multicolumn{6}{r}{Summa}&\multicolumn{2}{c|}{Thaler}&47324&5203&42366&4945
\end{tabular}
\end{center}
\noindent\hspace*{\parindent}Den Totalwerth der Witwenpensionen vergrössere ich um seinen sechsten Theil wegen der Waisenpensionen, nicht sowohl deswegen, weil \textsc{Kritter} in seinem Gutachten dieses Verhältniss angenommen hat, als weil dasselbe sehr nahe aus meiner eignen Discussion der Erfahrungen bei unsrer Witwenkasse hervorgeht [S. 154]. Sodann wird der Totalwerth der Beiträge abgezogen, wodurch sich der reine Werth der Verbindlichkeit der Kasse gegen die 42 verheiratheten Mitglieder ergibt. Für die 9 jetzt unverheiratheten wird endlich schlechthin pro rata zugesetzt. Diese Rechnungen stehen dann folgendermaassen:
\begin{center}
\footnotesize\setlength{\tabcolsep}{3pt}
\begin{tabular}{l r|r}
&Zinsfuss \(3\tfrac{1}{2}\) Proc.&Zinsfuss 4 Proc.\\
\shortstack[l]{Witwenpensionen der 42\\verheiratheten Mitglieder}&47324 Thl.&42366 Thl.\\
Waisenpensionen . . . . . . . . . . . . . . . . . . . . . .&7887 ,, &7061 ,, \\[5pt]
Pensionen . . . . . . . . . . . . . . . . . . . . . . . . . . .&55211 ,, &49427 ,, \\
Beiträge . . . . . . . . . . . . . . . . . . . . . . . . . . . .&5203 ,, &4945
\end{tabular}
\end{center}
